\section{Exact probabilistic termination and values}
\label{sec:markov}

\status{Finite rational-chain search, exact arithmetic replay, and closed-class
countercertificates are implemented in Python. The probability-space bridge and
its Lean formalization remain proposed.}

\subsection{The input semantics must include probability mass}
Let $S$ be a finite nonempty state set and let
$P\in\Q^{S\times S}$ satisfy
\[
 P_{st}\geq0,\qquad \sum_tP_{st}=1.
\]
The input specifies an initial state $s_0$, a terminal set $T\subseteq S$, a
nonnegative rational cost $c(s)$ per nonterminal visit, and a terminal payoff
$g(t)\in[0,1]\cap\Q$. Every terminal state is absorbing; its cost is zero.
The implementation stores $g(s)=0$ off $T$.

Write $X_n$ for the Markov chain and
$\tau=\inf\{n\geq0:X_n\in T\}$. The positive output means
\begin{equation}
 \Prob_{s_0}(\tau<\infty)=1,
 \quad h(s_0)=\E_{s_0}\tau,
 \quad v(s_0)=\E_{s_0}\sum_{k=0}^{\tau-1}c(X_k),
 \quad u(s_0)=\E_{s_0}g(X_\tau).
 \label{eq:mc-goals}
\end{equation}
If $g$ is the indicator of a successful terminal subset, $u$ is the success
probability. With arbitrary fractional terminal payoffs it is an expected
payoff, not automatically the probability of a specified event. The returned
field is therefore named \code{expected_terminal_payoff}; a regression test
uses payoff $1/3$ to preserve this distinction.

The validator rejects floating probabilities, negative entries, nonunit row
sums, nonabsorbing terminals, and nonzero terminal costs. There is no tolerance
parameter and no automatic normalization. These checks are part of the source
semantics. Multiplying a stochastic row by a normalizing factor after export
would change the problem unless that normalization were itself justified.

The worker's complete decision question is almost-sure absorption in $T$.
Exact costs and terminal values accompany positive answers. On a negative
answer it supplies a certified positive lower bound for nontermination, not
necessarily the exact nontermination probability or an exact total cost.
This scope matters when some nonterminating behavior has zero running cost.

\subsection{Why a harmonic equation is insufficient}
A tempting certificate for a probability is a vector $u$ satisfying $u=Pu$
with the intended boundary values. This is insufficient without a condition
that makes the boundary reachable in the appropriate probabilistic sense.
For the one-state nonterminal self-loop $P=[1]$, every scalar $u$ satisfies
$u=Pu$. No such scalar by itself proves that the process ever reaches a
terminal state; there is no terminal state at all.

Likewise, an exact solution of a floating-point-derived linear system is exact
only for the rounded model. A weighted-word identity concerns products along
finite words. Neither kind of evidence establishes that an infinite-horizon
limit has the claimed boundary values. The missing ingredient is a
\emph{properness witness}: a finite nonnegative potential that pays at least
one unit for every nonterminal step.

\subsection{The certificate and its local checker}
The certificate supplies a set $R\subseteq S$ containing $s_0$, closed under
all positive-probability successors, and three rational vectors $h,v,u$ on
$R$. The support need not be the smallest reachable set; it only has to be
forward closed. Outside $R$, the serialized vectors contain no values.

The checker verifies $h,v\geq0$, $0\leq u\leq1$, the terminal conditions
\[
 h(t)=v(t)=0,\qquad u(t)=g(t)\qquad(t\in R\cap T),
\]
and, for every $s\in R\setminus T$,
\begin{align}
 h(s)&=1+\sum_{t\in R}P_{st}h(t),\label{eq:h}\\
 v(s)&=c(s)+\sum_{t\in R}P_{st}v(t),\label{eq:v}\\
 u(s)&=\sum_{t\in R}P_{st}u(t).\label{eq:u}
\end{align}
All equalities are exact rational identities. The checker performs
matrix--vector products; it does not call Gaussian elimination, compute
strongly connected components, or trust a solver-supplied residual norm.

\begin{theorem}[Properness and exact values]
\label{thm:mc}
If the positive checker accepts these data, then \eqref{eq:mc-goals} holds.
Moreover the same claims hold from every state in $R$ with that state's
certificate values.
\end{theorem}
\begin{proof}
Forward closure ensures that a chain started in $R$ stays there. Extend the
nonterminal indicator by zero on $T$. Iterating \eqref{eq:h} for a finite
number $N$ of steps gives the exact identity
\begin{equation}
 h(s)=\sum_{k=0}^{N-1}\Prob_s(\tau>k)+\E_s[h(X_N)].
 \label{eq:finite-h}
\end{equation}
The sum is $\E_s[\min(\tau,N)]$. Its terms and the residual are nonnegative.
Letting $N$ increase shows $\E_s\tau\leq h(s)<\infty$, and hence
$\Prob_s(\tau=\infty)=0$. One can avoid this first limit by observing directly
that
\[
 (N+1)\Prob_s(\tau>N)
 \leq \sum_{k=0}^{N}\Prob_s(\tau>k)
 \leq h(s).
\]
Thus the survival probability tends to zero.

Since $R$ is finite, $H=\max_{t\in R}h(t)$ is finite. The residual in
\eqref{eq:finite-h} is at most $H\Prob_s(\tau>N)$, because $h$ vanishes
on terminals. It tends to zero. Therefore the finite sums converge exactly to
$h(s)$, proving the expected-time equality rather than only an upper bound.

Similarly, iterating \eqref{eq:v} yields
\[
 v(s)=\E_s\sum_{k=0}^{N-1}c(X_k)+\E_s[v(X_N)].
\]
The residual is at most $V\Prob_s(\tau>N)$, where
$V=\max_{t\in R}v(t)$. Terminal costs vanish and costs are nonnegative, so
monotone convergence identifies the limit with the total pretermination cost.
Its value is $v(s)$.

Finally, \eqref{eq:u} also holds at terminals, because they are absorbing.
Thus $u(s)=\E_s[u(X_N)]$ for every $N$. On $\{\tau\leq N\}$ this
random variable equals $g(X_\tau)$; on its complement it lies in $[0,1]$.
The complement's probability tends to zero, giving
$u(s)=\E_s[g(X_\tau)]$. No assertion that an arbitrary harmonic solution is
unique has been used without first proving properness.
\end{proof}

\subsection{A finite-horizon proof interface with explicit error bounds}
The preceding proof offers a particularly small first Lean target. Let
$\mu_N$ be the exact finite distribution after $N$ steps, defined recursively
by $\mu_0=\delta_s$ and $\mu_{N+1}=\mu_NP$. Put
\[
 t_N=\sum_{x\notin T}\mu_N(x),\quad
 C_N=\sum_{k=0}^{N-1}\sum_x\mu_k(x)c(x),\quad
 G_N=\sum_{x\in T}\mu_N(x)g(x).
\]
The checker theorem can first establish, entirely through finite sums and
induction on $N$,
\begin{align}
 t_N&\leq\frac{h(s)}{N+1},\label{eq:tail-bound}\\
 0\leq h(s)-\sum_{k=0}^{N-1}t_k
       &\leq\frac{Hh(s)}{N+1},\label{eq:time-bound}\\
 0\leq v(s)-C_N&\leq\frac{Vh(s)}{N+1},\label{eq:cost-bound}\\
 0\leq u(s)-G_N&\leq\frac{h(s)}{N+1}.\label{eq:payoff-bound}
\end{align}
These are rational inequalities for each natural $N$, not numerical
experiments. Terminal absorption makes $t_N$ nonincreasing; the exact finite
identities then give the bounds as in the proof. In particular,
$u(s)-G_N=\sum_{x\notin T}\mu_N(x)u(x)$, so its sign and bound are explicit.

A formalization can prove these finite-sum identities before introducing an
infinite path measure. It can then identify $t_N$ with the probability of a
hitting-time event using the standard probability semantics, or use the
sequences directly for a source language whose denotation is given as the
limit of finite execution distributions. The latter interface still needs a
source theorem; it must not call a finite-vector recurrence an operational
probability statement without that link.

The bounds are conservative. Geometric examples have exponentially decaying
tails, whereas \eqref{eq:tail-bound} is only inverse-linear. Its advantages are
that it follows from the same small certificate and makes every limiting step
explicit. A later worker can certify geometric drift or multi-step
contraction for sharper rates without changing the basic absorption checker.

\subsection{Search: graph structure first, rational solving second}
The algorithm first computes the states reachable from $s_0$ in the support
graph, whose edges are exactly the entries $P_{st}>0$. It decomposes this graph
into strongly connected components and looks for a reachable bottom component
containing no terminal state. If one exists, search returns the negative
certificate below. Unreachable traps are irrelevant to this query and are not
allowed to spoil a positive answer.

Otherwise let $U=R\setminus T$, let $Q=P_{U,U}$, and form
\[
 A=I-Q,
 \qquad b_h=\one,
 \qquad b_v=c|_U,
 \qquad b_u(s)=\sum_{t\in T}P_{st}g(t).
\]
The search solves the three right-hand sides simultaneously:
\begin{equation}
 A\,[h\ v\ u]=[b_h\ b_v\ b_u].
 \label{eq:matrix-solve}
\end{equation}
The prototype uses exact Gauss--Jordan elimination with rational pivots. It sets
terminal values from the original input and returns the full vectors on $R$.
The checker validates their signs and all original-model equations afresh.

Why is $A$ invertible in the positive case? Every reachable nonterminal state
has a support path to a terminal; otherwise its reachable subgraph would
contain a nonterminal bottom component. Choose one such path from each state.
Finiteness gives a common length bound $m$ and a positive minimum path
probability $\delta$. Conditional on not having terminated, the chain has
probability at least $\delta$ to terminate within the next $m$ steps. Hence
\[
 \Prob_s(\tau>km)\leq(1-\delta)^k.
\]
The transient matrix powers tend to zero, the series
$\sum_{k\geq0}Q^k$ converges, and it is the inverse of $I-Q$. This also shows
that absorption in a finite chain entails finite expected absorption time.
The corresponding assertion is not automatic for infinite-state chains.

A dense solve has $O(|U|^3)$ rational field-operation cost. Exact bit complexity
can be much larger because intermediate numerators and denominators grow.
The graph phase is linear in the explicit support graph, although the shipped
matrix representation also scans zero entries. Replay uses $O(|R|^2)$ exact
arithmetic operations in the dense representation and never factors or
inverts a matrix.

For larger instances, \code{fmpq_mat_solve} in FLINT is a concrete rational
linear-algebra backend \cite{flint}; modular or fraction-free methods can
replace the prototype solve without enlarging the checker. Storm is another
candidate search backend for richer probabilistic models \cite{storm}. Its
model export must preserve exact probabilities. In particular, Storm's
explicit transition-file interface is documented as using floating input and
not supporting the same exact mode as suitable symbolic-model inputs. Merely
passing \code{--exact} through a wrapper is not a proof that the original
probability model was retained.

\subsection{Negative certificate: a reachable closed nonterminal class}
A negative certificate is a nonempty set $C\subseteq S\setminus T$ and a
finite path $s_0,s_1,\ldots,s_k$ ending in $C$. The checker verifies
\begin{equation}
 P_{s_i s_{i+1}}>0\quad(i<k),\qquad
 P_{xy}>0\Longrightarrow y\in C\quad(x\in C).
 \label{eq:trap}
\end{equation}
It returns the exact rational
\[
 \delta=\prod_{i=0}^{k-1}P_{s_i s_{i+1}}>0.
\]
For a path of length zero, this product is one. The checker need not verify
that $C$ is strongly connected or minimal; closedness and nonterminality are
sufficient. Those stronger properties help the search find a small object,
but are not obligations of the certificate.

\begin{theorem}[Closed-class refutation]
\label{thm:trap}
An accepted negative certificate proves
$\Prob_{s_0}(\tau=\infty)\geq\delta>0$. Thus absorption is not almost sure.
\end{theorem}
\begin{proof}
The specified finite path has probability $\delta$. After it enters $C$, the
row-sum and closure conditions imply that all subsequent transitions stay in
$C$ with probability one. Since $C$ contains no terminal, the chain never
terminates on that event. A terminal cannot have appeared earlier on the path:
it is absorbing and could not be left to reach $C$.
\end{proof}

Every negative instance of the finite-chain almost-sure question has such a
certificate: some reachable bottom component is nonterminal. Together with
Theorem~\ref{thm:mc}, this makes the search a decision procedure on the stated
finite rational fragment, subject only to implementation resources. The
checker handles an intentionally simpler object than the SCC algorithm used
by search.

The negative verdict alone does not imply that total expected \emph{cost} is
infinite: a nonterminal closed class can have zero costs. It does imply that
expected hitting time is infinite, because $\tau=\infty$ has positive
probability. A source-level report must keep these two quantities separate.

\subsection{Worked examples and adversarial distinctions}
For a two-state chain with a retry state $r$ and successful terminal $t$, let
\[
 P=\begin{pmatrix}3/4&1/4\\0&1\end{pmatrix},
 \qquad c(r)=3/2,
 \qquad g(t)=1.
\]
The certificate is
\[
 h=(4,0),\qquad v=(6,0),\qquad u=(1,1).
\]
For example $4=1+(3/4)4$ and $6=3/2+(3/4)6$. The checker therefore proves
almost-sure termination, expected time four, expected cost six, and terminal
success probability one. The infinite sequence that remains at $r$ forever is
still a syntactically possible path, but it has probability zero. Replacing
the probabilistic choice by an antagonist that can choose the self-loop makes
universal eventual termination false. This example is shared with the game
demonstration precisely to expose the difference in quantifiers.

Now let a start state go to a successful terminal with probability $2/3$ and
to a nonterminal absorbing state with probability $1/3$. The singleton trap
and the one-edge path to it certify nontermination probability at least
$1/3$. No linear solve is needed. Changing that edge to probability zero must
invalidate the path certificate, even if the same vertices remain in the
serialized model.

A very small positive exit probability can also have a very large expected
waiting time. For exit probability $1/2^{40}$, the exact expected time is
$2^{40}$. A tolerance that rounds the exit probability to zero changes an
almost-surely terminating chain into a nonterminating one. Conversely,
rounding a truly zero edge upward can erase a trap. The exact input and exact
residual discipline address this semantic discontinuity; floating residual
size alone does not.

\subsection{Extensions that require additional certificates}
For a Markov decision process, graph SCCs must be replaced by the relevant
end-component analysis, and a policy becomes part of the certificate. A
properness equation for \emph{one} policy proves neither almost-sure
termination under every policy nor an optimum among policies. One concrete
sufficient certificate for a universal time bound is a potential satisfying
\[
 h(s)\geq1+\sum_tP(s,a,t)h(t)
\]
for every enabled action at every nonterminal state. Exact optimal values need
Bellman inequalities for all alternatives and an attaining policy, together
with the appropriate properness conditions. These are specific next steps,
not claims that the present chain solver handles nondeterministic scheduling.

Similarly, exact reachability probabilities in a chain that is not almost
surely absorbed in the original $T$ can be reduced by making nonterminal bottom
classes explicit failure terminals. That reduction preserves the reachability
event but not the original expected hitting time or total running cost. A
future adapter must state which quantity it preserves. The current negative
result deliberately returns only its proved lower bound rather than making an
unchecked reduction or overstating the available value information.
